% =====================================================================
%  VIII. CONCLUSION
%  Numbers repeat Section V and must stay identical to it.
% =====================================================================
\section{Conclusion}
\label{sec:conclusion}

A user who tries a configuration, judges the answer and resends the prompt
to a stronger one pays for every attempt, and none of the routing studies we
located measures that cost. We defined it as the escalation tax, split it
into escalation and stickiness, and measured it on three pools: 1{,}200
gradable prompts run through all eight cells of a Qwen3 size-by-mode grid on
a laptop GPU with metered board energy, thirteen public models priced by
their providers, and six hosted cells in which reasoning effort is a paid
setting.

The tax is large, and most of it comes from where a user stays rather than
how they climb. Leaving the 8B thinking cell selected pays a tax of 12.96
times the minimum sufficient energy and answers 79.8\% of prompts, while a
frugal climb pays 0.89 times and answers 92.2\%. The same ordering holds on
the public pool, with taxes of 7.96 against 2.69 times the minimum
sufficient cost, and on the hosted ladder, 4.35 against 0.22.

On the gradable families thinking was never the cheaper way to buy accuracy.
Below 8B, the next size up without thinking answered more prompts than
thinking at the same size, for 3.8 to 5.5 times less energy per solved
prompt. Open-ended requests, graded by a judge that agrees with a human only
moderately, differed: thinking cost about twice the energy per solved prompt
and the largest thinking cell was the most accurate, yet the frugal climb
still paid the least tax, 1.46 times the minimum sufficient energy against
5.39 for the sticky default. Routing
one tier too low cost 2.71 times as much energy as routing one tier too
high. Counted in tokens, that asymmetry reverses, and sticking with the
largest model looks cheaper than trying thinking first, the opposite of what
the meter shows. Cost has to be measured in the unit that is actually paid.
On held-out prompts no router beat the frugal climb. The best, logistic
regression on TF-IDF features, spent 1.30 times its energy per solved prompt
and answered 3.7 points fewer, while spending 64\% less than leaving the 8B
thinking cell selected and 37\% less than a self-verifying cascade, both per
solved prompt. Charging
routers for their own energy barely changed their ranking, since none cost
more than 7.6\,J per decision. The loss to the climb came from recovering
failed routes at the top of the grid.

For a user without a router, the evidence points to a short rule: start
small with reasoning off, move up a size before turning reasoning on, and
treat extended reasoning as a choice made per prompt rather than a setting
left on. A one-shot router did not improve on that rule when every failure
was recognised. It earned its place against the habits that start high, and
where failures go unnoticed.
